\documentclass[10pt,a4paper]{amsart}
\usepackage[margin=19mm]{geometry}
\usepackage{amsmath,amssymb,mathtools}
\usepackage[scr=rsfs,cal=euler]{mathalfa}
\usepackage{xfrac}
\usepackage[unicode,hidelinks,bookmarks=false]{hyperref}
\newcommand{\ie}{i.e.\ }
\newcommand{\cf}{cf.\ }
\newcommand{\resp}{respectively\ }
\newcommand{\Subsection}{\subsection}
\providecommand{\nobreakdash}{\nobreak\hbox{-}}
\setlength{\emergencystretch}{2em}
\multlinegap0pt
\usepackage{amssymb}
\usepackage{mathtools}
\usepackage{bm}
\newtheorem{lemma}{Lemma}
\newcommand{\op}{\mathrm{op}}
\title{\textbf{Distributional Limits for Eigenvalues of Graphon Kernel Matrices}}
\author{Behzad Aalipur}
\date{Source: arXiv:2601.04584v5 (CC BY 4.0)}
\begin{document}
\maketitle

\noindent\emph{Source and licence.} \textbf{Distributional Limits for Eigenvalues of Graphon Kernel Matrices}. Version arXiv:2601.04584v5 (CC BY 4.0), distributed by arXiv under the Creative Commons Attribution 4.0 International licence, http://creativecommons.org/licenses/by/4.0/, as recorded in the arXiv licence metadata for this exact version and shown on the abstract page https://arxiv.org/abs/2601.04584v5. Full licence notice: \texttt{licenses/arxiv-2601.04584-CC-BY-4.0.txt}.\par\smallskip

\noindent\textit{Editorial scope.} Counted unit: the article's lemma lem:expansion (source0.tex lines 1251--1270). The article's complete original proof of it is delivered with it (source0.tex lines 1272--1452, one \texttt{proof} environment of 181 lines). No mathematics is written, repaired or re-derived: the statement and the proof are the article's own lines, in the article's own notation, and no retained line is substituted or re-flowed.  No further source-local context is needed: every cross-reference of every retained line resolves inside the two delivered spans.  Nothing was omitted from any retained span, so the package carries no omission record.  No retained line contains a citation, so the wrapper carries no bibliography and no bibliography entry.  No retained line contains a figure environment, an image-inclusion command or drawing code, so no image is delivered and the package declares no asset dependency.  Editorial formatting only, in addition: an AMS \texttt{amsart} preamble that restates the article's own declarations and macros used by the retained lines, namely the environments \texttt{lemma} on separate counters, together with the standard packages those lines need.  The retained environments print in this document as Lemma 1 in order of appearance, whereas the article prints the same items with the numbering of its own full document.  The complete native source of the article as distributed in the arXiv e-print for this exact version is bundled unchanged as \texttt{sources/192302/source0.tex}.
\begin{lemma}[Second-order Rayleigh-Schrödinger expansion for a simple eigenvalue]
\label{lem:expansion}
Let $\mathbf{M} \in \mathbb{R}^{n \times n}$ be symmetric, and fix $r \in \{1,\dots,n\}$ such that 
$\lambda_r := \lambda_r(\mathbf{M})$ is simple with eigengap
\[
\gamma_r := \min_{k \neq r}|\lambda_r(\mathbf{M}) - \lambda_k(\mathbf{M})| > 0.
\]
Define the perturbed matrix $\widehat{\mathbf{M}} = \mathbf{M} + \mathbf{E}$, and let $\mathbf{u}_r$
be the unit eigenvector of $\mathbf{M}$ corresponding to $\lambda_r$.
If $\|\mathbf{E}\|_{\op} < \gamma_r/2$, then
\[
\widetilde{\lambda}_r - \lambda_r
=
\mathbf{u}_r^\top \mathbf{E}\mathbf{u}_r + R,
\]
where $\widetilde{\lambda}_r := \lambda_r(\widehat{\mathbf{M}})$ and the remainder satisfies
\[
|R| \le \frac{2\|\mathbf{E}\|_{\op}^2}{\gamma_r}.
\]
\end{lemma}

\begin{proof}
We follow the block decomposition with respect to the eigenspace of $\lambda_r$.
Let $\{\mathbf{u}_1,\dots,\mathbf{u}_n\}$ be an orthonormal eigenbasis of $\mathbf{M}$ with
$\mathbf{M}\mathbf{u}_k=\lambda_k(\mathbf{M})\mathbf{u}_k$ and $\mathbf{u}_r$ corresponding to $\lambda_r$.
Define
\[
\mathbf{Q}
:=
[\mathbf{u}_1,\dots,\mathbf{u}_{r-1},\mathbf{u}_{r+1},\dots,\mathbf{u}_n]
\in \mathbb{R}^{n\times(n-1)},
\]
so that $\mathbf{Q}$ spans the orthogonal complement of $\mathbf{u}_r$ and
\(
\mathbf{Q}^\top \mathbf{Q}=\mathbf{I}_{n-1},
\,
\mathbf{Q}\mathbf{Q}^\top=\mathbf{I}_n-\mathbf{u}_r\mathbf{u}_r^\top.
\)

Every $\mathbf{x}\in\mathbb{R}^n$ admits a unique decomposition
\[
\mathbf{x}=\mathbf{u}_r\xi+\mathbf{Q}\mathbf{z},
\qquad
\xi\in\mathbb{R},\ \mathbf{z}\in\mathbb{R}^{n-1}.
\]

Let $\widetilde{\mathbf{u}}_r$ be the unit eigenvector of $\widehat{\mathbf{M}}$
associated with $\widetilde{\lambda}_r$, and write
\[
\widetilde{\mathbf{u}}_r=\mathbf{u}_r\xi+\mathbf{Q}\mathbf{z},
\qquad
\xi^2+\|\mathbf{z}\|^2=1.
\]

The eigenvalue equation
\[
(\mathbf{M}+\mathbf{E})\widetilde{\mathbf{u}}_r
=
\widetilde{\lambda}_r \widetilde{\mathbf{u}}_r
\]
becomes
\[
(\mathbf{M}+\mathbf{E})(\mathbf{u}_r\xi+\mathbf{Q}\mathbf{z})
=
\widetilde{\lambda}_r(\mathbf{u}_r\xi+\mathbf{Q}\mathbf{z}).
\]

Projecting onto $\mathbf{u}_r$ and onto $\text{span}(\mathbf{Q})$ yields
\[
\begin{cases}
\lambda_r\xi + \mathbf{u}_r^\top \mathbf{E}(\mathbf{u}_r\xi+\mathbf{Q}\mathbf{z})
= \widetilde{\lambda}_r \xi, \\[0.3em]
\mathbf{Q}^\top \mathbf{M}\mathbf{Q}\mathbf{z}
+ \mathbf{Q}^\top \mathbf{E}(\mathbf{u}_r\xi+\mathbf{Q}\mathbf{z})
= \widetilde{\lambda}_r \mathbf{z}.
\end{cases}
\]

Let $\alpha:=\widetilde{\lambda}_r-\lambda_r$.
Rewrite the second equation as
\[
\bigl(\mathbf{L}_0 + \mathbf{\Delta}\bigr)\mathbf{z}
=
-\xi\,\mathbf{Q}^\top\mathbf{E}\mathbf{u}_r,
\]
where
\[
\mathbf{L}_0 := \mathbf{Q}^\top(\mathbf{M}-\lambda_r\mathbf{I})\mathbf{Q},
\qquad
\mathbf{\Delta} := \mathbf{Q}^\top(\mathbf{E}-\alpha\mathbf{I})\mathbf{Q}.
\]

By the eigengap assumption, $\|\mathbf{L}_0^{-1}\|_{\op}\le 1/\gamma_r$.
Under $\|\mathbf{E}\|_{\op}<\gamma_r/2$,
$\mathbf{L}_0+\mathbf{\Delta}$ is invertible and
\[
\|(\mathbf{L}_0+\mathbf{\Delta})^{-1}\|_{\op}
\le \frac{2}{\gamma_r}.
\]

Hence
\[
\mathbf{z}
=
-\xi\,(\mathbf{L}_0+\mathbf{\Delta})^{-1}
\mathbf{Q}^\top\mathbf{E}\mathbf{u}_r.
\]

From the first projected equation,
\[
\alpha\xi
=
\xi\,\mathbf{u}_r^\top\mathbf{E}\mathbf{u}_r
+
\mathbf{u}_r^\top\mathbf{E}\mathbf{Q}\mathbf{z}.
\]

If $\xi\neq0$, divide by $\xi$ and substitute $\mathbf{z}$:
\[
\alpha
=
\mathbf{u}_r^\top\mathbf{E}\mathbf{u}_r
-
\mathbf{u}_r^\top\mathbf{E}\mathbf{Q}
(\mathbf{L}_0+\mathbf{\Delta})^{-1}
\mathbf{Q}^\top\mathbf{E}\mathbf{u}_r.
\]

We now expand the inverse via a resolvent identity:
\[
(\mathbf{L}_0+\mathbf{\Delta})^{-1}
=
\mathbf{L}_0^{-1}
-
\mathbf{L}_0^{-1}\mathbf{\Delta}\mathbf{L}_0^{-1}
+
\mathbf{R}_3,
\]
where
\(
\|\mathbf{R}_3\|_{\op}
\le C \|\mathbf{E}\|_{\op}^2/\gamma_r^3.
\)

Substituting gives
\[
\alpha
=
\mathbf{u}_r^\top\mathbf{E}\mathbf{u}_r
-
\mathbf{u}_r^\top\mathbf{E}\mathbf{Q}
\mathbf{L}_0^{-1}
\mathbf{Q}^\top\mathbf{E}\mathbf{u}_r
+
R_3,
\]
with
\(
|R_3|
\le C\|\mathbf{E}\|_{\op}^3/\gamma_r^2.
\)

Since
\[
\mathbf{Q}\mathbf{L}_0^{-1}\mathbf{Q}^\top
=
\sum_{k\ne r}
\frac{1}{\lambda_k-\lambda_r}
\mathbf{u}_k\mathbf{u}_k^\top,
\]
we obtain the explicit second-order Rayleigh–Schrödinger term
\[
\mathbf{u}_r^\top\mathbf{E}\mathbf{Q}
\mathbf{L}_0^{-1}
\mathbf{Q}^\top\mathbf{E}\mathbf{u}_r
=
\sum_{k\ne r}
\frac{(\mathbf{u}_k^\top\mathbf{E}\mathbf{u}_r)^2}
{\lambda_r-\lambda_k}.
\]

Therefore,
\[
\widetilde{\lambda}_r-\lambda_r
=
\mathbf{u}_r^\top\mathbf{E}\mathbf{u}_r
+
\sum_{k\ne r}
\frac{(\mathbf{u}_k^\top\mathbf{E}\mathbf{u}_r)^2}
{\lambda_r-\lambda_k}
+
R_3,
\]
with
\[
|R_3|
\le
C\frac{\|\mathbf{E}\|_{\op}^3}{\gamma_r^2}.
\]

This yields the full second-order Rayleigh–Schrödinger expansion.
\end{proof}

\end{document}
